\chapter{Introduction}
\label{ch:intro}

\section{Contexte}
Dans un environnement produit toujours plus rapide, le \gls{po} joue un
rôle pivot entre la stratégie business, l'équipe de développement et la
voix de l'utilisateur. Ses tâches récurrentes --- rédaction de
\textit{user stories}, grooming d'épics, priorisation du backlog,
planification de sprint --- exigent à la fois rigueur méthodologique et
capacité d'arbitrage sous incertitude. La généralisation des modèles
de langage de grande taille (\gls{llm}) ouvre la voie à une nouvelle
classe d'outils d'aide à la décision capables d'\emph{augmenter}
l'expertise du PO, sans s'y substituer.

\section{Problématique}
Ce mémoire s'articule autour de la question suivante :
\begin{quote}\itshape
Dans quelle mesure une architecture \gls{rag} multi-modèles couplée
à des algorithmes de priorisation classiques peut-elle assister un
Product Owner dans ses décisions opérationnelles, tout en préservant
la traçabilité des recommandations et la maîtrise de l'humain~?
\end{quote}

\section{Objectifs}
\begin{enumerate}
  \item Concevoir une architecture logicielle modulaire séparant la
        couche RAG, la couche de priorisation et l'interface.
  \item Implémenter sept algorithmes de priorisation issus de la
        littérature Agile et les normaliser sur une échelle commune.
  \item Démontrer la valeur d'un routage multi-modèles
        (GPT-4o~/~Mistral Large) selon la nature de la tâche.
  \item Garantir la traçabilité par citation systématique des
        sources documentaires.
  \item Livrer un produit utilisable en démonstration réelle, déployé
        en production sur Vercel et Render.
\end{enumerate}

\section{Plan du mémoire}
Le chapitre~\ref{ch:sota} présente l'état de l'art du rôle de PO,
des frameworks Agile, des méthodes de priorisation et des
architectures RAG. Le chapitre~\ref{ch:agile} détaille la
méthodologie Scrum appliquée au projet. Les chapitres~\ref{ch:archi}
et~\ref{ch:impl} décrivent respectivement la conception et
l'implémentation. Le chapitre~\ref{ch:eval} présente l'évaluation
quantitative et qualitative de POSTIE. Le chapitre~\ref{ch:disc}
discute les limites observées. Enfin, le chapitre~\ref{ch:ccl}
conclut et propose des perspectives.
